\documentclass[11pt]{article}
\usepackage[margin=1in]{geometry}
\usepackage{amsmath, amssymb, amsthm}
\usepackage{booktabs}
\usepackage{array}
\usepackage{hyperref}

\hypersetup{colorlinks=true, linkcolor=blue, urlcolor=blue, citecolor=blue}

\newtheorem{definition}{Definition}
\newtheorem{assumption}{Assumption}
\newtheorem{theorem}{Theorem}
\newtheorem{corollary}{Corollary}
\newtheorem{remark}{Remark}

\title{Operationalizing the BP Decomposition on AutoResearch:\\Revised Theorem, Assumptions, and Validation Status}
\author{}
\date{April 13, 2026}

\begin{document}
\maketitle

\begin{abstract}
This note replaces the earlier AutoResearch BP draft with a narrower and more defensible formulation.
The original note was close to the Beneventano--Poggio (BP) framework, but it compressed three
non-trivial bridges into one sentence: the move from BP's model index to full architectures, the
move from fixed per-step cost to context-dependent average cost, and the move from noisy scalar
losses to a theorem-level verification object.

The strongest current theorem is best stated \emph{per axis}, not as a shared two-axis identity.
For a fixed axis $\alpha \in \{\mathrm{wall}, \mathrm{token}\}$, if AutoResearch admits an
architecture-indexed packed latent-mode family, if success is defined by a latent-loss threshold,
and if the exact BP-style cost factor is replaced in practice by an arithmetic mean cost
$\bar\kappa_\alpha$, then
\[
\Delta_\alpha(d \mid d_0)
=
\log \frac{\bar\kappa_\alpha(d_0)}{\bar\kappa_\alpha(d)}
+ \phi_\alpha(d_0,d)
+ G_\alpha(d)
- \epsilon_\alpha(d)
+ R_\alpha(d \mid d_0),
\]
where $R_\alpha$ is the Jensen remainder induced by replacing the exact effective cost with the
reported arithmetic mean.

This is weaker than the earlier two-axis statement, but it is the strongest defensible form on the
current evidence. The present repository evidence supports the usefulness of cost accounting and
strongly supports the need for repeated evaluation under noisy verification. It does \emph{not}
yet identify the full $\phi + G - \epsilon$ decomposition empirically, and therefore does not yet
justify a stronger theorem-estimation package.
\end{abstract}

\section{Introduction}

The motivating question is unchanged: how should one compare AutoResearch architectures for
iterative training research under deployment budgets? The intended answer comes from the BP
decomposition, which separates expected certified log-speedup into four channels:

\begin{enumerate}
\item per-step cost,
\item within-mode competence,
\item information generation,
\item routing mismatch.
\end{enumerate}

The earlier draft was directionally right but too fast in execution. It claimed that AutoResearch
inherits the BP decomposition ``directly'' after three modifications:

\begin{enumerate}
\item split certified time into wall-clock and cost axes,
\item replace the noisy verifier with a latent loss,
\item fold context pressure into per-step cost.
\end{enumerate}

After auditing the original note against \texttt{BP.pdf} and after reviewing the current repository
pilot, the right conclusion is stricter:

\begin{quote}
AutoResearch appears compatible with a BP-style decomposition under explicit additional
assumptions, but the strongest current theorem is a \emph{single-axis conditional reduction with
an explicit remainder term}, not a free two-axis corollary.
\end{quote}

This revised note makes that status explicit.

\section{Theorem-Level Objects}

\begin{definition}[Candidate space]
Fix a research substrate with immutable data preparation and evaluation harness.
Let $Y$ be the set of syntactically valid \texttt{train.py} files that compile and complete training
within the platform budget.
\end{definition}

\begin{definition}[Observed loss and latent loss]
For a candidate $y \in Y$, let $V(y,\omega) \in \mathbb{R}_{\ge 0}$ be the observed validation loss
returned by one stochastic training run under randomness $\omega$.
Define the latent loss
\[
L(y) := \mathbb{E}_\omega[V(y,\omega)].
\]
Lower values are better.
\end{definition}

\begin{definition}[Theorem-level success event]
Fix a target loss threshold $q^\star$.
The theorem-level success event is
\[
\mathbf{1}\{L(y)\le q^\star\}.
\]
That is, the theorem is written for threshold success on \emph{latent} loss, not for one-shot noisy
evaluations. Noisy evaluations enter only at the estimation layer.
\end{definition}

\begin{definition}[Architecture]
An AutoResearch architecture is a design tuple
\[
d = (M,\pi,\Omega,W,K,p),
\]
where $M$ is the base LLM, $\pi$ is the orchestration/routing policy, $\Omega$ is the tool set,
$W$ is external memory or workspace structure, $K$ is the active context budget per LLM call,
and $p$ is the number of parallel agents.
\end{definition}

\begin{definition}[Axis-specific certified time]
For an axis $\alpha$ and target threshold $q^\star$, let $T_\alpha^{(d)}(q^\star,\delta)$ be the
certified time required by architecture $d$ to reach latent-loss success with probability at least
$1-\delta$ when measured on axis $\alpha$. The expected certified log-speedup relative to baseline
$d_0$ is
\[
\Delta_\alpha(d \mid d_0)
:=
\log T_\alpha^{(d_0)}(q^\star,\delta)
- \log T_\alpha^{(d)}(q^\star,\delta).
\]
\end{definition}

\section{Where The Earlier Draft Was Too Fast}

The earlier draft relied on BP correctly at the algebraic level, but it compressed three extra
bridges that are not automatic:

\begin{enumerate}
\item BP Theorem 21 is indexed by a model-like object. AutoResearch varies a full architecture.
\item BP uses a fixed per-step cost $\kappa(M)$. AutoResearch uses context-dependent cost
      $\kappa_\alpha(d,\Xi)$ and then reports arithmetic means.
\item The deployed verifier is noisy. One-shot observed loss is not the right theorem-level object.
\end{enumerate}

None of these issues kills the framework. But each requires either an explicit assumption or an
explicit remainder term.

\section{Revised Assumptions}

\begin{assumption}[Architecture-indexed packed family]
There exists a latent mode variable $S \in \{1,\dots,m\}$ for the AutoResearch task family such
that, for each fixed axis $\alpha$ and architecture $d$, the pair $(S,d)$ defines a BP-style
packed family with:
\begin{enumerate}
\item a mode posterior $\pi_{\alpha,D}^d$,
\item a routing allocation $q_{\alpha,D}^d$,
\item a within-mode certified scale $t_{0,\alpha}(d,s)$.
\end{enumerate}
This is the bridge from BP's model index to an architecture index.
\end{assumption}

\begin{assumption}[Latent-loss threshold regularity]
The latent loss $L(y)$ is measurable and finite for all candidates of interest, and the event
$\mathbf{1}\{L(y)\le q^\star\}$ is the theorem-level verification object.
Observed noisy evaluations are unbiased but noisy measurements of this latent object and are not
used directly inside the theorem statement.
\end{assumption}

\begin{assumption}[Axis-specific effective cost and reported mean cost]
For each axis $\alpha$ and architecture $d$, there exists an exact effective BP-style cost factor
$\tilde\kappa_\alpha(d) > 0$ such that the axis-specific certified time admits the BP log-time
factorization
\[
T_{\alpha,q}^{(d)}(s)
=
\tilde\kappa_\alpha(d)\, c_{\mathrm{sim}}\, \frac{t_{0,\alpha}(d,s)}{q(s)}.
\]
In practice the experiment reports the arithmetic mean cost
\[
\bar\kappa_\alpha(d) := \mathbb{E}[\kappa_\alpha(d,\Xi)],
\]
where $\Xi$ denotes the run/path dependence induced by context pressure and solver dynamics.
Define the Jensen gap
\[
J_\alpha(d) := \log \bar\kappa_\alpha(d) - \log \tilde\kappa_\alpha(d).
\]
Assume $J_\alpha(d)$ is finite.
\end{assumption}

\begin{remark}
If $\tilde\kappa_\alpha(d)$ is taken as the geometric mean effective cost, then $J_\alpha(d)\ge 0$
by Jensen's inequality. Under small within-architecture cost variance one expects
\[
J_\alpha(d)=O\!\left(\frac{\mathrm{Var}[\kappa_\alpha(d,\Xi)]}{\mathbb{E}[\kappa_\alpha(d,\Xi)]^2}\right),
\]
which is the quantity the protocol should estimate.
\end{remark}

\begin{assumption}[Prior-matched baseline]
There exists a baseline architecture $d_0$ whose routing is prior-matched for the packed family
used in the BP reduction.
\end{assumption}

\section{Strongest Current Theorem}

\begin{theorem}[Conditional single-axis decomposition with Jensen remainder]
Fix an axis $\alpha \in \{\mathrm{wall}, \mathrm{token}\}$.
Under Assumptions 1--4, for any architecture $d$ compared against a prior-matched baseline $d_0$,
the expected certified log-speedup satisfies
\[
\Delta_\alpha(d \mid d_0)
=
\log \frac{\bar\kappa_\alpha(d_0)}{\bar\kappa_\alpha(d)}
+ \phi_\alpha(d_0,d)
+ G_\alpha(d)
- \epsilon_\alpha(d)
+ R_\alpha(d \mid d_0),
\]
where
\[
\phi_\alpha(d_0,d)
:=
\mathbb{E}_{S \sim \pi_\alpha}\!\left[
\log \frac{t_{0,\alpha}(d_0,S)}{t_{0,\alpha}(d,S)}
\right],
\]
\[
G_\alpha(d)
:=
I_{\pi_\alpha}(S : D_\alpha^d),
\]
\[
\epsilon_\alpha(d)
:=
\mathbb{E}_{D_\alpha}\!\left[
\mathrm{KL}(\pi_{\alpha,D}^d \,\|\, q_{\alpha,D}^d)
\right],
\]
and
\[
R_\alpha(d \mid d_0)
:=
J_\alpha(d)-J_\alpha(d_0).
\]
Equivalently, if the exact effective cost $\tilde\kappa_\alpha$ is reported instead of the arithmetic
mean $\bar\kappa_\alpha$, the remainder vanishes and the BP four-term identity holds exactly on
that axis.
\end{theorem}

\begin{proof}
Fix an axis $\alpha$.
By Assumption~3, the axis-specific certified time has the multiplicative form required by BP
Theorem~21, with the architecture $d$ playing the role of BP's effective model index and with
$\tilde\kappa_\alpha(d)$ playing the role of BP's per-step cost:
\[
\log T_{\alpha,q}^{(d)}(s)
=
\log \tilde\kappa_\alpha(d)
+
\log c_{\mathrm{sim}}
+
\log t_{0,\alpha}(d,s)
-
\log q(s).
\]
Applying BP Theorem~21 on this axis yields
\[
\Delta_\alpha(d \mid d_0)
=
\log \frac{\tilde\kappa_\alpha(d_0)}{\tilde\kappa_\alpha(d)}
+ \phi_\alpha(d_0,d)
+ G_\alpha(d)
- \epsilon_\alpha(d).
\]
By the definition of the Jensen gap,
\[
\log \tilde\kappa_\alpha(d)
=
\log \bar\kappa_\alpha(d) - J_\alpha(d),
\]
so substituting for both $d_0$ and $d$ gives
\[
\log \frac{\tilde\kappa_\alpha(d_0)}{\tilde\kappa_\alpha(d)}
=
\log \frac{\bar\kappa_\alpha(d_0)}{\bar\kappa_\alpha(d)}
+ J_\alpha(d)-J_\alpha(d_0).
\]
This is exactly the stated decomposition with remainder
$R_\alpha(d \mid d_0)=J_\alpha(d)-J_\alpha(d_0)$.
\end{proof}

\begin{remark}[Why the theorem is single-axis]
The earlier revised note assumed that $\phi$, $G$, and $\epsilon$ are shared across wall-clock and
cost axes. That may still be a good approximation, but it is not needed for the theorem and is not
yet verified empirically. The stronger shared-axes statement should therefore be treated as a
\emph{future corollary} that requires additional evidence, not as part of the base theorem.
\end{remark}

\begin{corollary}[Shared-axis corollary, conditional]
If future work shows that the same mode partition, posterior family, and routing allocations are
valid on both axes and that the within-mode certified scales are axis-invariant up to pricing, then
\[
\phi_{\mathrm{wall}}=\phi_{\mathrm{token}},\qquad
G_{\mathrm{wall}}=G_{\mathrm{token}},\qquad
\epsilon_{\mathrm{wall}}=\epsilon_{\mathrm{token}},
\]
and the two-axis statement from the earlier draft follows by applying the single-axis theorem to
each axis separately.
\end{corollary}

\section{Claim Status Table}

\begin{table}[h]
\centering
\begin{tabular}{p{0.31\linewidth} p{0.2\linewidth} p{0.39\linewidth}}
\toprule
Item & Status & Meaning in this note \\
\midrule
BP four-term algebra for packed families & Proved algebra & Inherited from BP Theorem 21 once the architecture-indexed packed family exists. \\
\addlinespace
Architecture index replacing BP model index & Structural assumption & Stated explicitly as Assumption 1. \\
\addlinespace
Threshold success on latent loss & Structural assumption & The theorem is written on latent-loss threshold success, not one-shot noisy evaluations. \\
\addlinespace
Arithmetic mean cost replacing exact effective cost & Structural assumption plus remainder & Captured by Assumption 3 and the Jensen remainder $R_\alpha$. \\
\addlinespace
Axis sharing of $\phi$, $G$, $\epsilon$ & Pending corollary & Not assumed in the base theorem; requires extra evidence. \\
\addlinespace
Estimators recover $\phi$, $G$, $\epsilon$ & Estimation assumption & Not part of the theorem; current implementation does not satisfy it. \\
\addlinespace
H1--H5 style pilot predictions & Empirical hypotheses & Must be tested experimentally; they are not theorem-level consequences. \\
\bottomrule
\end{tabular}
\caption{What is proved, assumed, or still empirical in the current formulation.}
\end{table}

\section{Empirical Consequences}

The theorem does not by itself prove the pilot hypotheses; it organizes them.
For the 2x2 design the most relevant empirical questions remain:

\begin{enumerate}
\item whether parallelism improves wall-clock certified time enough to offset coordination costs;
\item whether external memory lowers effective cost by relieving context pressure;
\item whether shared memory lowers routing mismatch relative to independent parallel agents;
\item whether the sign of the parallelism contrast flips across axes;
\item whether higher context pressure measurably enlarges the Jensen gap and effective cost.
\end{enumerate}

These are empirical hypotheses, not theorem-level corollaries.

\section{Validation Status On The Current Repository}

The repository already contains a 12-run CPU pilot and a follow-up noise assay.
The present evidence should be summarized as follows.

\begin{table}[h]
\centering
\begin{tabular}{p{0.33\linewidth} p{0.57\linewidth}}
\toprule
Question & Status on current repository evidence \\
\midrule
Is the BP reduction structurally credible? &
Yes, once the architecture, latent-loss, and cost assumptions are stated explicitly. \\
\addlinespace
Was the full four-term decomposition identified? &
No. In the current pilot, $\phi$ is hardcoded to zero and the implemented $G,\epsilon$ estimators are degenerate. \\
\addlinespace
Was the stronger two-axis shared-term statement validated? &
No. The current evidence does not justify axis-sharing for $\phi$, $G$, or $\epsilon$. \\
\addlinespace
Was the latent-loss layer operationally necessary? &
Yes. Repeated evaluations show large noise and strong best-of-many optimism. \\
\addlinespace
Was the original cost axis reproduced literally? &
No. The CPU substrate uses token-only cost and drops GPU pricing by design. \\
\bottomrule
\end{tabular}
\caption{Validation status of the current repository evidence.}
\end{table}

\subsection{Noise assay}

To test whether the latent-loss layer is genuinely needed, the baseline training configuration was
repeated five times and the pilot-selected best $d_{10}$ candidate was repeated five times.

\begin{center}
\begin{tabular}{lccccc}
\toprule
Config & Mean val\_bpb & Std & Range & Pilot best & Replicated? \\
\midrule
Baseline & 0.8246 & 0.0359 & 0.0930 & --- & --- \\
Best $d_{10}$ candidate & 0.8694 & 0.0501 & 0.1198 & 0.7554 & No \\
\bottomrule
\end{tabular}
\end{center}

The pilot-selected best $d_{10}$ commit did \emph{not} replicate:
its repeated-evaluation mean was about $0.114$ worse than the pilot-reported best value.
This is consistent with strong selection bias under noisy one-shot evaluation.

\subsection{Deterministic Phase 02 calibration update}

After the original pilot-and-noise pass, the repository was upgraded with a deterministic
calibration workflow on the CPU CIFAR-10 substrate. The completed Phase~02 comparison was the
single-agent control $d_{00}$ versus the single-agent external-memory architecture $d_{10}$.

\begin{center}
\begin{tabular}{lcc}
\toprule
Metric & $d_{00}$ & $d_{10}$ \\
\midrule
Best-of-rep mean val\_bpb & 0.8905 & 0.9151 \\
All-runs mean val\_bpb & 1.0026 & 1.0127 \\
Total training runs & 47 & 69 \\
Mean runs / replicate & 9.4 & 13.8 \\
Modes with $\ge 2$ accepted edits & 1 & 1 \\
\bottomrule
\end{tabular}
\end{center}

These results sharpen the empirical status in an important way.

\begin{enumerate}
\item On the descriptive best-of-replicate view, $d_{10}$ still does not outperform $d_{00}$.
      If anything, the difference points in the wrong direction for the simple memory-helps story.
\item On the stricter all-runs view with current post-hoc labeling, the quality effect is
      negligible.
\item $d_{10}$ does appear to improve iteration throughput: it produces more runs per replicate.
\item The accepted-mode family remains too thin to count as empirical identification of the full
      $\phi + G - \epsilon$ package.
\end{enumerate}

So the deterministic calibration supports a more refined reading:
\emph{memory may help the cost / throughput channel without helping the quality channel}.
That is perfectly compatible with the revised theorem, which separates cost accounting from the
other terms, but it does \emph{not} count as validation of the full four-term decomposition.

The repository therefore contains a genuine decision split:
the descriptive best-of-replicate workflow summary justified an exploratory move to the parallel
$d_{01}$ / $d_{11}$ cells, while the stricter accepted-mode gate still points toward
structured search as the cleaner next step.
Those later parallel runs should be read as exploratory override evidence, not as a calibration
gate that has been cleanly passed.

\section{What Is Validated vs. Pending}

\subsection*{Validated in this revision}

\begin{itemize}
\item The strongest current theorem is a conditional single-axis BP reduction with an explicit
      Jensen remainder.
\item The empirical setting requires a latent-loss/repeated-evaluation layer.
\item The current pilot provides some signal on cost-channel differences across architectures.
\end{itemize}

\subsection*{Still pending}

\begin{itemize}
\item a non-degenerate operational mode definition with full labeling coverage;
\item correct estimators for $\phi_\alpha$, $G_\alpha$, and $\epsilon_\alpha$;
\item direct estimation of the Jensen gap through within-architecture cost variance;
\item a disciplined repeated-evaluation protocol integrated into architecture comparison;
\item a genuine context-pressure sweep that enters higher $c/K$ regimes.
\end{itemize}

\section{Conclusion}

The earlier AutoResearch note was directionally right but too compressed to count as a rigorous
theorem statement. This revision narrows the claim deliberately:

\begin{enumerate}
\item the theorem is written on threshold success for latent loss;
\item the strongest theorem is single-axis, not two-axis shared by default;
\item replacing exact effective cost by arithmetic mean cost incurs an explicit remainder;
\item the estimator layer is separate from the theorem and is not yet validated.
\end{enumerate}

The correct research posture is therefore:

\begin{quote}
keep the decomposition, narrow the theorem, repair the estimators and protocol, and only then ask
whether the full empirical decomposition is identifiable.
\end{quote}

\section*{References}

\begin{enumerate}
\item P. Beneventano and T. Poggio. \emph{A Mathematical Framework for Building Agent Systems:
Task-First Design, Equivalence, and Model Choice under Deployment Budgets}. Under review, 2026.
\item Original AutoResearch note in this repository: \texttt{autoresearch\_bp.pdf}.
\item Repository validation bundle: \texttt{theory\_validation\_bp\_20260412/}.
\end{enumerate}

\end{document}
